\documentclass[12pt,a4paper]{article}
\usepackage[utf8]{inputenc}
\usepackage[T1]{fontenc}
\usepackage[polish]{babel}
\usepackage{amsmath}
\usepackage{amsfonts}
\usepackage{amssymb}
\usepackage{booktabs}
\usepackage{geometry}
\usepackage{enumitem}
\usepackage{parskip}
\geometry{margin=2cm}

\title{ZTBSK (2026L), 4. Sieciowe usługi nazw}
\author{Piotr Gugnowski, wydział EiTI, nr indeksu 320690}
\date{25 Maj 2026}

\begin{document}
\maketitle

\section*{Pytania i odpowiedzi}

% -----------------------------------------------------------------------
\subsection*{1. Do czego służą sieciowe usługi nazw?}

Sieciowe usługi nazw służą do translacji symbolicznych nazw (zrozumiałych dla
człowieka) na adresy sieciowe (zrozumiałe dla urządzeń sieciowych) i odwrotnie.
Każdy program korzystający z łączności sieciowej, który potrafi posługiwać się
nazwami, musi być klientem takiej usługi; w praktyce jest on konsolidowany
z dynamiczną biblioteką zwaną resolverem. Dzięki temu użytkownicy mogą odwoływać
się do zasobów sieciowych przy użyciu czytelnych nazw (np.\ \texttt{www.example.com})
zamiast numerycznych adresów IP.

Poza tłumaczeniem nazw na adresy, usługi te mogą przechowywać szerszą gamę
informacji katalogowych (np.\ dane o użytkownikach, zasobach organizacji,
uprawnieniach), co stanowi podstawę dla usług katalogowych takich jak LDAP.

% -----------------------------------------------------------------------
\subsection*{2. Jakie usługi zaliczają się do sieciowych usług nazw?}

Do sieciowych usług nazw zalicza się przede wszystkim:

\begin{description}
  \item[DNS (\emph{Domain Name System})]
        Rozproszona, hierarchiczna baza danych obejmująca swoim zasięgiem całą
        sieć Internet. Służy do tłumaczenia w pełni kwalifikowanych nazw domenowych
        (FQDN) na adresy IP i odwrotnie, a także do przechowywania innych informacji
        o zasobach sieciowych (rekordy MX, NS, TXT i inne).
  \item[LDAP (\emph{Lightweight Directory Access Protocol})]
        Nowoczesna usługa katalogowa o zasięgu lokalnym (instytucjonalnym).
        W przeciwieństwie do DNS, nie istnieje jedna ogólnoświatowa struktura łącząca
        wszystkie serwery LDAP; każda instytucja tworzy własne lokalne serwery.
        Najważniejszym zastosowaniem LDAP jest autoryzacja użytkowników,
        przechowywanie danych kontaktowych i zarządzanie kontami w obrębie organizacji.
\end{description}

% -----------------------------------------------------------------------
\subsection*{3. Co to jest DNS?}

DNS (\emph{Domain Name System}) jest to \textbf{rozproszona, hierarchiczna baza
danych}. Wszelkie informacje przechowywane są w postaci \textbf{rekordów}
składających się z pięciu elementów:
\begin{itemize}
  \item \textbf{nazwa:} identyfikator pozwalający na odnalezienie rekordu w bazie,
  \item \textbf{czas życia} (TTL, \emph{Time To Live}): wyrażony w sekundach
        okres, przez który dozwolone jest przechowywanie rekordu w lokalnych buforach,
  \item \textbf{klasa rekordu:} identyfikator tworzący odrębną przestrzeń nazw;
        jedyną powszechnie stosowaną klasą jest \texttt{IN} (Internet),
  \item \textbf{typ rekordu:} identyfikator mówiący o rodzaju wartości (np.\
        \texttt{A} dla adresu IPv4, \texttt{AAAA} dla IPv6, \texttt{CNAME} dla
        nazwy kanonicznej, \texttt{MX} dla serwera pocztowego),
  \item \textbf{wartość:} istotna informacja o obiekcie.
\end{itemize}

Rekordy pogrupowane są w tzw.\ \textbf{strefy} (\emph{domains}), tworzące strukturę
drzewiastą. Istnieje strefa główna (\emph{root domain}), w której zarejestrowane są
nazwy stref najwyższego poziomu (TLD, \emph{Top Level Domains}), a w każdej strefie
można rejestrować kolejne strefy niższego poziomu.

System jest \textbf{rozproszony}, ponieważ za przechowywanie informacji o rekordach
każdej strefy odpowiedzialna jest inna grupa serwerów.

% -----------------------------------------------------------------------
\subsection*{4. Jakie są typy serwerów systemu DNS?}

Ze względu na pełnioną funkcję serwery DNS dzielą się na dwa główne typy:

\begin{description}
  \item[Serwery rekursywne (buforujące)]
        Służą do rozwiązywania zapytań na potrzeby lokalnych klientów. Gdy klient
        chce poznać adres IP dla danej nazwy, wysyła zapytanie do lokalnego serwera
        rekursywnego, który przejmuje cały ciężar poszukiwania odpowiedzi. Przy braku
        informacji w pamięci podręcznej (cache) serwer sam odpytuje kolejne serwery
        w hierarchii DNS (poczynając od serwerów strefy głównej), aż do uzyskania
        odpowiedzi autorytatywnej. Wyniki są buforowane, by uniknąć powtarzania
        zapytań.
  \item[Serwery autorytatywne]
        Udzielają odpowiedzi autorytatywnych dotyczących stref, za które są
        odpowiedzialne. Dzielą się dalej na:
        \begin{itemize}
          \item \textbf{serwer główny} (\emph{master/primary}), przechowujący
                oryginał danych danej strefy,
          \item \textbf{serwery pomocnicze} (\emph{slave/secondary}), przechowujące
                kopię danych strefy, automatycznie aktualizowaną z serwera głównego.
        \end{itemize}
\end{description}

Technicznie jeden serwer DNS może pełnić jednocześnie kilka funkcji: być serwerem
autorytatywnym dla jednej lub kilku stref i jednocześnie obsługiwać jako serwer
rekursywny grupę lokalnych maszyn. Ze względów bezpieczeństwa (ataki DoS) zaleca
się jednak rozdzielanie tych funkcji.

% -----------------------------------------------------------------------
\subsection*{5. Które serwery, i kiedy, wysyłają odpowiedzi autorytatywne?}

Do udzielania odpowiedzi autorytatywnych uprawnione są \textbf{wyłącznie serwery
odpowiedzialne za strefę}, w której zarejestrowana jest dana nazwa, tj.\ serwery
autorytatywne (główny i pomocnicze) dla tej konkretnej strefy.

Odpowiedź autorytatywna jest oznaczona specjalną flagą w pakiecie DNS i wysyłana
jest zawsze, niezależnie od tego, czy jest pozytywna (rekord istnieje) czy
negatywna (rekord nie istnieje). Uzyskanie odpowiedzi autorytatywnej kończy
rozwiązywanie nazwy i powoduje odesłanie odpowiedzi do klienta.

Serwery rekursywne nie wysyłają odpowiedzi autorytatywnych; mogą jedynie
przekazywać informacje ze swojej pamięci podręcznej (które zostały wcześniej
uzyskane od serwera autorytatywnego) lub inicjować dalsze zapytania w imieniu
klienta.

% -----------------------------------------------------------------------
\subsection*{6. Co oznacza skrót FQDN?}

FQDN to skrót od angielskiego \emph{Fully Qualified Domain Name} (pol.\
\textbf{w pełni kwalifikowana nazwa domenowa}). Aby dotrzeć do informacji zawartej
w konkretnym rekordzie DNS, konieczne jest podanie listy nazw stref począwszy od
strefy najwyższego poziomu aż do nazwy rekordu. Taki pełny identyfikator, zawierający
kompletną ścieżkę przez hierarchię DNS, nazywany jest właśnie FQDN.

Przykładowo \texttt{www.example.com.} jest FQDN, gdzie kropka na końcu oznacza
strefę główną (\emph{root domain}). W codziennym użyciu kropka na końcu jest
pomijana, lecz w konfiguracji serwerów DNS ma ona istotne znaczenie, gdyż odróżnia
FQDN od nazwy względnej.

% -----------------------------------------------------------------------
\subsection*{7. Które z serwerów DNS muszą być dostępne dla wszystkich komputerów
w sieci?}

Dla wszystkich komputerów w sieci \textbf{lokalnej} musi być dostępny
\textbf{lokalny serwer rekursywny (buforujący)}. To on przyjmuje zapytania od
klientów i realizuje ich rozwiązanie. Adres tego serwera jest konfigurowany
w ustawieniach sieciowych każdego komputera (ręcznie lub przez DHCP).

Serwery autorytatywne dla stref publicznych (w szczególności \textbf{serwery
strefy głównej} i serwery TLD) muszą być dostępne dla serwerów rekursywnych
z całego Internetu, ponieważ każde rozwiązywanie nazwy ostatecznie opiera się
na hierarchii zapoczątkowanej przez te serwery.

Serwery autorytatywne dla stref prywatnych (lokalne strefy organizacji) wystarczy
udostępnić w obrębie sieci lokalnej.

% -----------------------------------------------------------------------
\subsection*{8. W jaki sposób serwery pomocnicze (\emph{slave}) dowiadują się,
że uległa zmianie zawartość strefy na serwerze głównym (\emph{master})?}

Stosowane są dwie komplementarne metody:

\begin{description}
  \item[Cykliczne odpytywanie serwera głównego (\emph{polling})]
        Serwery pomocnicze w regularnych odstępach czasu pytają serwer główny,
        czy dane strefy uległy zmianie (porównując numer seryjny w rekordzie SOA).
        Jeśli numer seryjny wzrósł, pobierają aktualizację (tzw.\ transfer strefy).
        Metoda ta zapewnia niezawodność, gdyż aktualizacja nastąpi zawsze, niezależnie
        od stanu połączenia w chwili zmiany.
  \item[Powiadomienia od serwera głównego (\emph{NOTIFY})]
        Serwer główny aktywnie informuje serwery pomocnicze o tym, że dane uległy
        zmianie. Zapewnia to szybką propagację zmian bez oczekiwania na kolejny
        cykl odpytywania.
\end{description}

W praktyce często łączy się obie metody, aby zapewnić jednocześnie niezawodność
(odpytywanie) i szybkość propagacji zmian (powiadomienia).

% -----------------------------------------------------------------------
\subsection*{9. Co to jest DNS dynamiczny? Jakie są z nim związane problemy?}

DNS był projektowany przy założeniu, że adresy IP są związane na stałe z komputerami,
a ewentualne zmiany (dodanie, usunięcie, zmiana adresu) są stosunkowo rzadkie, stąd
typowe czasy życia rekordów (TTL) są rzędu tygodnia. \textbf{Dynamiczny DNS} to próba
wykorzystania systemu DNS do obsługi dynamicznie zmieniających się przydziałów
adresów IP (np.\ w środowiskach korzystających z DHCP).

\textbf{Problemy związane z dynamicznym DNS:}

\begin{itemize}
  \item \textbf{Wyścigi i hazardy przy niskich wartościach TTL.}
        Pozornym rozwiązaniem jest drastyczne skrócenie TTL rekordów. Jednak gdy
        czasy życia stają się bliskie czasom potrzebnym na uzyskanie odpowiedzi od
        serwera, możliwe stają się wyścigi prowadzące do uzyskania nieaktualnej
        informacji.
  \item \textbf{Niemożność buforowania, wzrost obciążenia serwerów.}
        Krótkie TTL uniemożliwiają efektywne buforowanie zapytań. Powoduje to
        generowanie znacznie większej liczby zapytań do serwerów, co obciąża je
        i zmniejsza ich wydajność.
  \item \textbf{Problem autoryzacji aktualizacji.}
        Przydziały dynamiczne są dokonywane poza serwerem DNS (np.\ przez serwer
        DHCP lub same stacje robocze), co oznacza, że serwer DNS musi akceptować
        żądania zmiany zawartości strefy przychodzące z zewnątrz. Szczególnie
        trudne jest zapewnienie bezpieczeństwa, gdy żądania przysyłają poszczególne
        stacje robocze; wielokrotnie wykrywano błędy umożliwiające hakerom zdalne
        manipulowanie zawartością serwera.
\end{itemize}

% -----------------------------------------------------------------------
\subsection*{10. Co to jest DNS podzielony? Jakie jest jego przeznaczenie?}

\textbf{DNS podzielony} (ang.\ \emph{split DNS}) to konfiguracja, w której serwery
DNS udzielają \textbf{różnych odpowiedzi w zależności od tego, kto zadaje pytanie}.
W tym celu definiuje się tzw.\ \textbf{widoki} (ang.\ \emph{views}) i określa,
którzy klienci mają otrzymywać dane z którego widoku:

\begin{itemize}
  \item zapytania z sieci lokalnej: ujawniane są wszystkie nazwy i adresy
        prywatne (pełny widok wewnętrzny),
  \item zapytania z zewnątrz: ujawniane są wyłącznie nazwy i adresy publiczne
        (widok zewnętrzny).
\end{itemize}

\textbf{Przeznaczenie DNS podzielonego} wynika z powszechnego stosowania prywatnych
adresów IP w sieciach lokalnych. Z jednej strony wygodnie jest posługiwać się
nazwami wewnątrz sieci lokalnej, z drugiej strony prywatne adresy nie są dostępne
spoza niej. Informacja o nich byłaby bezużyteczna dla zewnętrznych użytkowników,
a jednocześnie potencjalnie przydatna dla intruzów. DNS podzielony rozwiązuje
ten problem, odseparowując wewnętrzną i zewnętrzną przestrzeń nazw przy użyciu
tych samych serwerów.

\end{document}
